% TOP_PROBLEM: {"id": 3259, "title": "Non-edges vs. feedback edge sets in digraphs", "category_id": 3, "category": {"id": 3, "name": "graph_theory", "display_name": "Graph Theory", "description": "Problems involving graphs, networks, and their properties.", "slug": "graph-theory", "order_index": 3, "created_at": "2026-07-31T15:26:25.671Z"}, "status": "open", "problem_number": "OPG-1793", "set_id": 12}
% FIRST_POSTED: 2026-10-01
% ENTRY_KIND: statement_only
% UnsolvedMath ID 3259; source catalogue code OPG-1793
% !TeX program = lualatex
% Definition and sources only; this file is not an LLM solution attempt.
\documentclass[11pt,a4paper]{article}
\usepackage[margin=25mm]{geometry}
\usepackage{fontspec}
\setmainfont{lmroman10-regular.otf}[BoldFont=lmroman10-bold.otf,
 ItalicFont=lmroman10-italic.otf,BoldItalicFont=lmroman10-bolditalic.otf]
\usepackage{amsmath,amssymb,mathtools}
\usepackage{unicode-math}
\setmathfont{latinmodern-math.otf}
\AtBeginDocument{\let\setminus\smallsetminus}
\usepackage{xurl,hyperref}
\hypersetup{colorlinks=true,urlcolor=blue}
\setcounter{secnumdepth}{0}
\setlength{\parindent}{0pt}
\setlength{\parskip}{5pt}
\setlength{\emergencystretch}{3em}
\begin{document}
\section*{Non-edges vs. feedback edge sets in digraphs}\label{3259}
{\small UnsolvedMath ID 3259 · OPG-1793 · Graph Theory · OpenGarden}
\subsection{Definitions and mathematical statement}
Let $D=(V,A)$ be a finite simple digraph with no directed cycle of length
one, two, or three. Consequently there are no loops or pairs of opposite
arcs. A non-edge is an unordered pair of distinct vertices with no arc
in either direction. Let
\[
 \gamma(D)=\bigl|\{\{u,v\}\subseteq V:uv\notin A, vu\notin A\}\bigr|.
\]
The unordered convention counts each missing adjacency once.

A feedback arc set is a subset $F\subseteq A$ such that $D-F$ has no
directed cycle. Equivalently, $F$ meets the arc set of every directed
cycle of $D$. Define
$\beta(D)=\min\{|F|:F\subseteq A,\ D-F\text{ is acyclic}\}$.
All deletions here are arc deletions, not vertex deletions.

The Chudnovsky--Seymour--Sullivan conjecture asks whether every such
3-free digraph satisfies
\[
                        \beta(D)\le\frac{\gamma(D)}2.
\]
Since $\beta(D)$ is an integer, the right-hand side can equivalently be
replaced by $\lfloor\gamma(D)/2\rfloor$. The condition allows undirected
triangles with a transitive orientation, and permits directed cycles
of length four or more.

One interpretation is that an almost-complete orientation with no short
directed cycle must be close to acyclic. More precisely, a linear ordering
of $V$ has a backward arc when its head precedes its tail. The minimum
number of backward arcs over all such orderings equals $\beta(D)$,
by taking a topological ordering after deleting a minimum feedback
arc set. This gives an equivalent ordering formulation of the target.
\subsection{Short English statement}
In a digraph without directed cycles of length at most three, can all directed cycles be destroyed by deleting at most half as many arcs as there are missing vertex pairs?
\subsection{Sources}
\begin{itemize}
\item[S1] ULAM.AI and the UnsolvedMath Contributors, supplied problems.json, record 3259 (OPG-1793), OpenGarden collection. Catalogue identity and source transcription; CC BY 4.0.
\url{https://huggingface.co/datasets/ulamai/UnsolvedMath}.
\item[S2] Open Problem Garden, Non-edges vs. feedback edge sets in digraphs. Original problem and discussion; the mathematical formulation below is expanded from the supplied source record.
\url{http://www.openproblemgarden.org/op/non_edges_vs_feedback_edge_sets_in_digraphs}.
\end{itemize}
\end{document}
